\documentclass[exercice,a4]{thedocument}%
\usepackage[french]{babel}%
\usepackage{tabularx}%
\usepackage{diagbox}%
\usepackage{array}%
\startdatakeys
    \source{21GENMATG11}
    \BO{2026}
    \cycle{4}
    \competencesDuSocle{CA,RA}
    \objectifsApprentissage{C4C1bc,C4C2bd,C4B2eb}
    \motsCles{probabilité}
\enddatakeys
\begin{document}%

\textbf{Partie 1}

\medskip
Dans cette première partie, on lance un dé bien équilibré à six faces numérotées de 1 à 6, puis on note le numéro de la face du dessus.

	\begin{enumerate}
		\item Donner sans justification les issues possibles.

		\item Quelle est la probabilité de l'évènement A : \og On obtient 2 \fg{} ?

		\item Quelle est la probabilité de l'évènement B : \og On obtient un nombre impair \fg{} ?
	\end{enumerate}

	\textbf{Partie 2}

\medskip

Dans cette deuxième partie, on lance simultanément deux dés bien équilibrés à six faces, un rouge et un vert. On appelle \og score \fg{} la somme des numéros obtenus sur chaque dé.

	\begin{enumerate}
		\item Quelle est la probabilité de l'évènement C : \og le score est 13 \fg{} ? Comment appelle-t-on un tel évènement ?

		\item Dans le tableau à double entrée donné en ANNEXE, on remplit chaque case avec la somme des numéros obtenus sur chaque dé.

		\begin{enumerate}
			\item Compléter, sans justifier, le tableau donné en ANNEXE à rendre avec la copie.

			\item Donner la liste des scores possibles.
		\end{enumerate}

		\item \begin{enumerate}
			\item Déterminer la probabilité de l'évènement D : \og le score est 10 \fg.

			\item Déterminer la probabilité de l'évènement E : \og le score est un multiple de 4 \fg.

			\item Démontrer que le score obtenu a autant de chance d'être un nombre premier qu'un nombre strictement plus grand que 7.
		\end{enumerate}
	\end{enumerate}

\vspace{0,5cm}

\begin{center}
	\begin{tabularx}{0.8\linewidth}{|>{\centering \arraybackslash}p{3cm}|*{6}{>{\centering \arraybackslash \rule[-3mm]{0mm}{10mm}}X|}}  \hline
	\backslashbox{\rule{0mm}{6mm}Dé rouge}{Dé vert}&1&2&3&4&5&6\\ \hline
	1& & & & & & \\ \hline
	2& & & & & & \\ \hline
	3& & & &7& & \\ \hline
	4& &6& & & & \\ \hline
	5& & & & & & \\ \hline
	6& & & & & & \\ \hline
\end{tabularx}
\end{center}

\vspace{5mm}

\showthesource
\end{document}
